\begin{textsample}{POS Dim 8 – human – Score 2.00 – rj/rj\_ef\_intro\_6\_p3.txt}  \label{ex:f8_pos_015}
O que se estima nesta transição de infância para adolescência, sem perder a continuidade processual do conhecimento?
O que considerar como \textbf{eixo} entre etapas, idade e interesse dos alunos quando inseridos no Ensino Fundamental Anos Finais, sem depreciar a capacidade de interlocução com o mundo de forma mais autônoma e crítica?
Para tal, reforça -se a importância de fortalecer a autonomia, oferecendo -lhes condições e ferramentas para que sejam críticos, reflexivos e que contribuam efetivamente para a mudança da \textbf{realidade} partindo de suas vivências, experiências e relação no cotidiano com múltiplos conhecimentos e fontes de informação.

% matched lemmas: eixo, realidade
\end{textsample}
